\section*{Chapter \ref{ch:iv2:etfs-futures-and-index-arbitrage} --- ETFs, Futures and Index Arbitrage}
\markboth{\protect\omTocMark{}{Answers: ETFs and Index Arbitrage}}{\protect\omTocMark{}{Answers: ETFs and Index Arbitrage}}

\begin{solution}{iq:iv2:etfs-futures-and-index-arbitrage:1}
Net asset value\index{net asset value} is \$40.00 per share, and the premium is 25 basis points.

The holdings are worth $50 \times 20 + 30 \times 50 + 10 \times 100 + 500 = 1\,000 + 1\,500 + 1\,000 + 500 = 4\,000$, for 100 ETF shares. And $0.10/40 = 0.25\%$.

\iqwrong{Leaving out the cash. It belongs to the shareholders, and it is often a large part of the difference for a fund near a dividend date.}
\end{solution}

\begin{solution}{iq:iv2:etfs-futures-and-index-arbitrage:2}
There is no arbitrage between \$39.96 and \$40.06. At a bid of \$40.09, you sell the ETF, buy the basket and create, and you make 3 cents a share.

To create you pay $\$40.00 + \$0.06$, so a sale above \$40.06 is a profit. To redeem you buy the ETF, then sell the basket for \$40.00 less \$0.04, so a purchase below \$39.96 is a profit. So the ETF's fair range is wider than the fund's own value, and it is set by creation and redemption\index{creation and redemption} costs.

\iqwrong{Using the same cost on both sides, or comparing the bid with the ask. You sell at the bid. And creating and redeeming have their own, possibly very different, costs.}
\end{solution}

\begin{solution}{iq:iv2:etfs-futures-and-index-arbitrage:3}
The fair price is $100 \times 1.02 - 1.01 = \$100.99$. At \$102, sell the forward and buy the stock with borrowed money.

Borrow \$100 for six months to buy the stock, which costs \$102 at the end. The \$1 dividend, in three months, earns $1\% $ for the remaining quarter, so it is worth \$1.01 at the end. You deliver the stock at \$102, repay \$102, and keep the \$1.01. So your locked-in profit is \$1.01, the amount by which \$102 exceeds \$100.99.

\iqwrong{Using $S(1 + rT)$ and ignoring the dividend. A holder of the forward gets no dividend, so the forward is cheaper by what the dividend is worth at delivery.}
\end{solution}

\begin{solution}{iq:iv2:etfs-futures-and-index-arbitrage:4}
One contract's notional is \$200\,000. You need 50 contracts for a beta of 1, and 60 for a beta of 1.2.

$4\,000 \times \$50 = \$200\,000$. Then $\$10\,000\,000/\$200\,000 = 50$. A beta of 1.2 means the portfolio moves like \$12 million of the index, which is 60 contracts.

\iqwrong{Hedging notional when the beta differs from 1. With 50 contracts, a beta of 1.2 leaves \$2 million of market exposure.}
\end{solution}

\begin{solution}{iq:iv2:etfs-futures-and-index-arbitrage:5}
The implied financing rate is 3\% a year. At a borrowing rate of 2.5\%, you buy the basket with borrowed money and sell the future.

The future's premium is $20/4000 = 0.5\%$ for a quarter, or 2\% a year, and $r - q = 2\%$, so $r = 3\%$. This is the cash and carry\index{cash and carry}: you are lending your money at the implied 3\% and borrowing at 2.5\%, a pick-up of 0.5\% a year before costs.

\iqwrong{Comparing the premium's 2\% with 2.5\%, and concluding the future is cheap. The holder of the basket also gets the dividend, so you have to add it back.}
\end{solution}

\begin{solution}{iq:iv2:etfs-futures-and-index-arbitrage:6}
The future implies 20 index points of dividends, a yield of 2\% a year.

Fair value without dividends is $4000 \times (1 + 5\% \times \tfrac14) = 4050$. The future at 4\,030 is 20 points lower, so that is the implied dividend\index{implied dividend}. And $20/4000 = 0.5\%$ for a quarter.

\iqwrong{Treating every gap between spot and futures as dividends. The gap also holds any financing spread, and in stress it can hold plenty of that.}
\end{solution}

\begin{solution}{iq:iv2:etfs-futures-and-index-arbitrage:7}
The index returns $-0.04\%$, the fund returns $-0.36\%$, and three times the index's return is $-0.12\%$.

The index is $1.02 \times 0.98 = 0.9996$. The fund is $1.06 \times 0.94 = 0.9964$. Each day's multiple is exact, but over two days, the return compounds differently, and it is lower by the extra volatility.

\iqwrong{Thinking a three-times fund gives three times the return over any horizon. It promises the daily multiple only.}
\end{solution}

\begin{solution}{iq:iv2:etfs-futures-and-index-arbitrage:8}
The fund loses about 1.82\%. Combined with the long index for those two days, you lose about 0.91\% on the \$2 you put in.

Up 10\% means the fund loses 10\%. Then the index falls $10/110 \approx 9.09\%$, and the fund gains 9.09\%. So $0.9 \times 1.0909 \approx 0.9818$. The index holding is flat, so with \$1 in each, you end with \$1.9818.

The one-day hedge is exact, but it stops being exact as soon as the two positions drift, since the fund's exposure is reset daily and you have not done the same to your position.

\iqwrong{Treating the inverse fund as a long-horizon hedge. To keep it a hedge, you have to rebalance between the two, as often as the fund resets.}
\end{solution}

\begin{solution}{iq:iv2:etfs-futures-and-index-arbitrage:9}
Sell \$10 million of the winner and buy \$10 million of the loser, a one-way turnover of 5\%. A cap-weighted fund needs no trades, because its weights move with the prices by themselves.

The holdings become \$110 million and \$90 million. Equal weight needs \$100 million of each. The total is \$200 million, so \$10 million is 5\%.

So cap-weighting has low turnover, and equal-weighting systematically sells winners and buys losers, and pays the trading cost to do so.

\iqwrong{Assuming that a cap-weighted fund must trade after a price move. It trades on changes in the index's membership or share counts.}
\end{solution}

\begin{solution}{iq:iv2:etfs-futures-and-index-arbitrage:10}
Your loss is \$50\,000, and \$50\,000 moves in cash today. The initial margin is \$100\,000 at the old price and \$97\,500 at the new one, so you must have \$47\,500 more than the \$50\,000 left in the account.

Your position's notional is $10 \times 50 \times 4000 = \$2\,000\,000$, and 5\% of that is \$100\,000. A 100-point fall on 10 contracts at \$50 is \$50\,000. So after the settlement, the account has \$50\,000, and the new requirement of $5\% \times \$1\,950\,000 = \$97\,500$ needs \$47\,500 more.

A forward with the same exposure would have no cash moving until settlement. Its gain or loss is carried as a credit risk. Futures turn market risk into a liquidity risk, and you need the cash today.

\iqwrong{Thinking a hedged position cannot cause trouble. A futures hedge against an illiquid long position can be margin-called in cash while the long cannot be sold.}
\end{solution}

\begin{solution}{iq:iv2:etfs-futures-and-index-arbitrage:11}
The fund loses about 2\% a month from the roll, for a compounded loss of about 21\% a year on the unchanged curve, before interest on its cash. An unchanged spot does not give an unchanged result because the fund keeps buying a higher-priced future that then falls toward spot.

Each month, it sells at 70 and buys at 71.40, so it holds $70/71.40 = 0.9804$ as many contracts. In a year, that is $0.9804^{12} \approx 0.788$, a loss of 21.2\%.

This curve is in contango\index{contango}: the future is above spot. It pays for storage and financing. The holder of a contract gets none of the benefit of having oil, and so the cost is charged to that holder.

\iqwrong{Saying the fund loses on the trade at each roll because it ``pays 1.40''. The roll itself is at market. The loss comes as the higher-priced contract falls to the front price over the month.}
\end{solution}

\begin{solution}{iq:iv2:etfs-futures-and-index-arbitrage:12}
Before the reset, assets are \$1.03 billion and exposure is \$3.03 billion. So it must buy \$60 million at the close. The flow is $L(L - 1)rA$. A minus-three fund of the same size buys \$120 million.

Its exposure was \$3 billion, and with 1\% up it is \$3.03 billion. Its assets rose by \$30 million. The target is $3 \times 1.03 = \$3.09$ billion, so it buys \$60 million.

In general, exposure after the move is $LA(1 + r)$, assets are $A(1 + Lr)$, and the target is $LA(1 + Lr)$. The difference is $LAr(L - 1)$.

At $L = -3$, $L(L - 1) = 12$, so it trades $12 \times 1\% \times \$1$ billion: it must also buy when the market rises. As a check at $L = 1$: nothing to do, and that is right.

\iqwrong{Thinking a fund that is short the market sells when the market rises. The index's rise has increased its short exposure and reduced its assets, so it must buy.}
\end{solution}

\begin{solution}{iq:iv2:etfs-futures-and-index-arbitrage:13}
Over a time $T$, the fund's value is $(S_T/S_0)^L\exp\left(-\tfrac12L(L - 1)\sigma^2T\right)$. With 30\% volatility and a flat index, the three-times fund returns about $-23.7\%$, and the minus-one fund about $-8.6\%$.

The log return of a daily-reset fund is about $L$ times the index's log return, with the volatility drag $-\tfrac12L^2\sigma^2T$ of the fund against $-\tfrac12\sigma^2T$ on the index. So the difference is $-\tfrac12(L^2 - L)\sigma^2T$.

For $L = 3$: $e^{-3 \times 0.09} = e^{-0.27} \approx 0.763$. For $L = -1$: $L(L - 1) = 2$, so $e^{-0.09} \approx 0.914$.

\iqwrong{Calling this a fee or a mistake in the fund. It is the arithmetic of rebalancing, and the fund does exactly what it says.}
\iqfollow{And if the index has trended?}{A trend can make the fund beat $L$ times the return: if the index rises 50\% with little volatility, the three-times fund is up nearly $1.5^3 - 1 = 237\%$. It depends on path.}
\end{solution}

\begin{solution}{iq:iv2:etfs-futures-and-index-arbitrage:14}
The new divisor is $130/60 \approx 2.167$. After a 10\% rise, the index is 66, and it would have been 66 without the split. So the split has no effect on the level, but it changes the weights.

The split takes the first stock to \$50, so the sum of prices is $50 + 50 + 30 = 130$. The index must stay at 60, so the divisor is $130/60$. If each stock rises 10\%, the sum is 143, and $143/2.167 = 66$. Without a split, the sum is 198, and $198/3 = 66$.

If only the first stock rose 10\% after the split, the index gains $5/2.167 \approx 2.31$, while before the split it would have gained $10/3 \approx 3.33$. So the split cut the stock's weight.

\iqwrong{Thinking the split leaves the index's sensitivity the same. A price-weighted index's weights are the share prices, so a split reduces that stock's importance.}
\end{solution}

\begin{solution}{iq:iv2:etfs-futures-and-index-arbitrage:15}
The index falls 0.8 points. A total-return version does not fall. The fair value of the future does not move on the day, because the dividend was already in the forecast.

$4000 \times 2\% \times 1\% = 0.8$. A price index does not reinvest, so it drops by the dividend. A total-return index adds the dividend back, so it is unchanged. The future's fair value is spot, less dividends to come. Spot falls by 0.8, and the dividends yet to come fall by 0.8, so the two cancel, up to interest.

\iqwrong{Calling a fall on the ex-date a loss. A holder gets the cash, so the total position is unchanged.}
\end{solution}

\begin{solution}{iq:iv2:etfs-futures-and-index-arbitrage:16}
You gain 20 points. The other side is whoever holds the future long; often that is a structured-product dealer laying off the long dividend exposure its book carries. For the exposure alone, you would trade a dividend future\index{dividend future} or a dividend swap.

You receive $2.5\%$ of 4\,000 instead of 2\%, which is 100 points rather than 80. The price of the future did not change, so the 20 points is your profit.

Dealers who sell structured products to investors often get long dividends as a by-product, and they sell them on. So implied dividends can sit below the expected level for long periods. The interest on the dividend has been ignored here.

\iqwrong{Treating a basis trade as risk-free. Your forecast of dividends is one of its main risks, and it matters more at long maturity.}
\end{solution}

\begin{solution}{iq:iv2:etfs-futures-and-index-arbitrage:17}
Over one year, the future costs 19 basis points and the ETF 8. They are equal at one month, so below that the future is cheaper.

The future costs $19H$ for a holding period of $H$ years. The ETF costs $7H + 1$. Setting them equal gives $12H = 1$. So the fixed spread of the ETF is what matters on a short horizon, and the running cost on a long one.

\iqwrong{Comparing only the fees. Implied financing is hidden in the futures price, and it moves. In stress it can rise far above 15 basis points.}
\end{solution}

\begin{solution}{iq:iv2:etfs-futures-and-index-arbitrage:18}
Fair value for the spread is about 30 points, so 4\,050 is fair. At 4\,060, the later future is about 10 points rich. Then, if you must stay long, you could buy spot through the basket at expiry, or take the near one to expiry and enter another way, since rolling costs an extra 10 points, about 0.25\%.

The spread is $4020 \times (5\% - 2\%) \times \tfrac14 \approx 30.15$. A calendar spread is the market price of financing for that quarter.

If you hold the near future, you can sell the spread (sell near, buy far) at 40 points, paying 10 over fair. Instead you could look for other exposure: hold to expiry, then buy the basket or an ETF, with its own costs. If those are under 10 points, you skip the roll.

\iqwrong{Treating the roll as a free event at the month's end. Everyone rolls in the same few days, so the calendar spread can carry a premium.}
\end{solution}

\begin{solution}{iq:iv2:etfs-futures-and-index-arbitrage:19}
The expected gap is $-0.16\%$ a year. Tracking error is the 0.05\% variation.

The fee costs 0.20\%, and lending earns 0.08\%. A cash holding of 1\% in a year when the index beats cash by 4\% costs 0.04\%. So $-0.20 + 0.08 - 0.04 = -0.16$.

Tracking difference\index{tracking difference} is the average gap, which is what you actually lose. Tracking error is its volatility. A fund can have low tracking error and a poor tracking difference.

\iqwrong{Using the fee as the whole cost of holding a fund. Lending can more than pay the fee, and the drag from cash reverses in a falling market.}
\end{solution}

\begin{solution}{iq:iv2:etfs-futures-and-index-arbitrage:20}
The no-arbitrage band is 4\,040 to 4\,064.

Above the band, you sell the future and buy the basket with borrowed money: $4000 \times (1 + (5\% - 2\%) \times \tfrac12) = 4060$, and with costs of 4 points, that gives 4\,064. Below the band, you buy the future and short the basket, and you lend the proceeds at 4.5\%, with a borrow fee of 0.3\% to pay. So $4000 \times (1 + (4.5\% - 0.3\% - 2\%) \times \tfrac12) = 4044$, and with the costs, 4\,040.

It is not symmetric because the two sides face different rates and only the short side pays a fee. Its width, 24 points, has two parts: 8 from the costs, and 16 from the gap in financing rates. So a firm with better financing has a tighter band and can arbitrage where you cannot.

\iqwrong{Thinking there is one fair value for a future. There is one for each cost of funding, so a future can sit between many firms' bands.}
\end{solution}

\begin{solution}{iq:iv2:etfs-futures-and-index-arbitrage:21}
You would make 0.30 on the ETF share if it all happened at once. But selling takes three days, with a standard deviation of about 0.87, so it is not arbitrage. The 95 is the better estimate. The published 100 is stale.

Buying at 95 and receiving bonds worth 95.50, less 0.20, leaves 0.30. The risk is $0.5\% \times \sqrt 3 \approx 0.87\%$. So one standard deviation is almost three times the profit.

Bonds trade rarely. The net asset value is made from stale prices or models, while the ETF trades now. In stress, the ETF is often the price discovery venue, and the ``discount'' is largely the true loss not yet in the marks. An ETF's discount is at most a modest arbitrage.

\iqwrong{Saying the ETF ``broke'' because of a discount. The discount reflects bonds that could not be sold at their marks.}
\end{solution}

\begin{solution}{iq:iv2:etfs-futures-and-index-arbitrage:22}
You net about 0.8 cents a share, or \$800 on 100\,000 shares. A basis move of 1 basis point against you is 0.4 cents a share, or \$400.

You earn 1 cent against fair value on the fill. The hedge costs half the future's spread: $0.5 \text{ bp} \times \$40 = 0.2$ cents. So $1 - 0.2 = 0.8$.

The risk is the gap between the ETF and the future: the future may not match the ETF's basket, and the future's value includes financing and dividend assumptions. So you are holding an extra risk until you offset it. Half your profit can go on a basis move of 1 basis point.

\iqwrong{Treating the hedge as perfect. A market maker's profit lies largely in how small this risk stays relative to the half-spread.}
\end{solution}

\begin{solution}{iq:iv2:etfs-futures-and-index-arbitrage:23}
The market maker has only the net \$10 million to hedge in the basket, which costs about \$2\,000 in total. Per dollar of ETF traded, on \$190 million, that is about 0.1 of a basis point. So an ETF's spread can be narrower than the basket's, because most trades are offset by other investors.

$\$10$ million at 2 basis points is \$2\,000. Over \$190 million, that is about 0.105 basis points. The offsetting trades never go to the basket, so only a little is needed to cover that part. The rest of the spread pays for risk.

\iqwrong{Assuming that an ETF cannot be cheaper than what it holds. Its liquidity is more than the sum of the parts as long as trading is two-way. In a one-way market, the spread goes back toward the basket's.}
\end{solution}

\begin{solution}{iq:iv2:etfs-futures-and-index-arbitrage:24}
The swap is cheaper by 10 basis points a year: cash costs funding plus 60, and the swap costs funding plus 50.

In cash you pay 30 and lose $15\% \times 2\% = 30$ on dividends. With the swap, you pay 40 and lose $5\% \times 2\% = 10$.

With a swap, you take on the dealer's credit risk, usually reduced by daily margin, and the dealer may end the trade. You have no vote, and the terms can change. The dealer hedges by holding the stock, and it may get better tax treatment or financing than you, and that is the source of the gap. Check the legal position, since tax rules on such swaps vary by country.

\iqwrong{Comparing only the two financing spreads. The dividend leak is the same size as the financing here, and it decides the answer.}
\end{solution}

\begin{solution}{iq:iv2:etfs-futures-and-index-arbitrage:25}
They must buy \$1.6 billion. That is 16 days at 100\% of volume, and 80 days at 20\%. The rule gives a price pressure of 8\%. You should doubt it, because 16 days of volume is far outside where the rule was measured.

Free float is $\$10$ billion $\times 80\% = \$8$ billion, and 20\% of that is \$1.6 billion. Then $1.6/0.1 = 16$, and $16/0.2 = 80$. The rule says $2\% \times \sqrt{16} = 8\%$.

Reasons to doubt it: this demand is announced, and that draws in sellers who would not trade otherwise, as well as arbitrageurs who hold stock in advance. Volume on the effective day is many times normal. So you may see only part of this price pressure, and part of that will reverse.

\iqwrong{Saying the funds must buy over 80 days. They want the closing price on the effective date, so they will buy nearly all of it at one moment.}
\end{solution}

\begin{solution}{iq:iv2:etfs-futures-and-index-arbitrage:26}
Your expected profit is zero. Your median profit is about \$0.41, and you lose money about 28\% of the time. A loss of about \$2 or more comes about once in 20 years. The decay is real for a typical path, and it is paid for with a rare, very large loss.

Each fund's value is $\prod_i(1 \pm 3r_i)$, and with independent, mean-zero daily returns, its expected value is exactly 1. So short positions in both have zero expected profit, despite each fund's strong decay. By the formula, with drag $\exp(-\tfrac12L(L - 1)\sigma^2T)$ and $x$ the index's multiple, the total value of the two funds is about $e^{-0.27}x^3 + e^{-0.54}x^{-3}$. At $x = 1$, that is about 1.35, so the profit is about 0.65 if the index ends flat. The most you can make, if realised volatility is exactly 30\%, is about 0.67.

You lose when the index's log return is above about $+28\%$ or below about $-37\%$. The value then grows with the cube of the index's gain, and that loss has no upper limit. A simulation of daily returns gives the median, the 28\% and the tail.

\iqwrong{Treating the decay as a yield. It is a typical-path result, and it has the same logic as selling a strangle: mean zero, and a positive median. Even that holds before the borrow fee, which is usually high on these funds.}
\end{solution}

\begin{solution}{iq:iv2:etfs-futures-and-index-arbitrage:27}
Each group sells \$2.4 billion, \$4.8 billion in all, which is 24\% of the last 15 minutes' liquidity. It looks like a short gamma hedge because the amount is proportional to the move and it trades with the move.

By the formula $L(L - 1)rA$. For $L = 3$: $6 \times (-4\%) \times \$10$ billion $= -\$2.4$ billion. For $L = -3$: $12 \times (-4\%) \times \$5$ billion $= -\$2.4$ billion.

Check the first by exposures: it was \$30 billion and is now \$28.8 billion, while the assets are $\$10 \times 0.88 = \$8.8$ billion, so the target is \$26.4 billion. That is a sale of \$2.4 billion. Check the second: it was $-\$15$ billion and is now $-\$14.4$ billion, the assets are $\$5 \times 1.12 = \$5.6$ billion, and the target is $-\$16.8$ billion. That is a sale of \$2.4 billion.

Someone short gamma has a delta that grows against them, and has to trade with the market. Here, the funds are short gamma in effect: their position per move is $L(L - 1)A$. So a late-day trend can be helped along by the funds' own flow.

\iqwrong{Thinking long and inverse funds cancel. They both trade with the move, so their flows add.}
\end{solution}

\begin{solution}{iq:iv2:etfs-futures-and-index-arbitrage:28}
The up-state probability is $5/9$. The forward price is about 108.35, and the futures price is about 108.75. They differ because the stock is positively correlated with interest rates. If the rate is not random, or not correlated with the stock, the difference is zero.

The stock must grow at the rate on average: $100 \times 1.04 = 120p + 84(1 - p)$, which gives $p = 20/36 = 5/9$.

The two-period zero-coupon bond is worth $\tfrac{1}{1.04}\left(\tfrac59/1.06 + \tfrac49/1.02\right) \approx 0.92292$, so the forward is $100/0.92292 \approx 108.35$. The futures price is the expected final price under this measure: $\tfrac59 \times 120 \times 1.06 + \tfrac49 \times 84 \times 1.02 \approx 108.75$.

The long future gets daily gains in the up state, when the rate is high and they can be invested at it, and losses in the down state, when they can be financed cheaply. So the holder has an advantage over a forward, and the futures price must be higher to offset it. This is the convexity adjustment\index{convexity adjustment}.

\iqwrong{Treating forwards and futures as the same at all times. They are close for short equity contracts, but they differ materially for long-dated contracts and for rates themselves.}
\end{solution}

\begin{solution}{iq:iv2:etfs-futures-and-index-arbitrage:29}
The weights are 0.54 on $A$ and 0.24 on $B$, per dollar of ETF. The remaining volatility is about 8.8\%, a variance reduction of 76\%.

The covariances among the hedge stocks are $0.04$, $0.03$ and $0.09$. The ETF's covariances with them are $0.8 \times 0.18 \times 0.2 = 0.0288$ and $0.7 \times 0.18 \times 0.3 = 0.0378$. Solving $\Sigma w = c$ gives $w = (0.54, 0.24)$. The explained variance is $w^\top c = 0.024624$ out of $0.0324$. What is left is $18\% \times \sqrt{0.24} \approx 8.8\%$.

More names give a lower variance in the sample, but your costs and the error in the estimates both grow. A large covariance matrix fits noise, so out of sample, weights can be unstable, and the real variance is higher than the fit says. So choose a small number of stable hedges.

\iqwrong{Using the two stocks' market-cap weights in the ETF. A minimum-variance hedge is not the same as a slice of the holdings.}
\end{solution}

\begin{solution}{iq:iv2:etfs-futures-and-index-arbitrage:30}
If you stay short, you make 6. With free flipping, you can expect 7, so the option is worth 1 point. With 0.5 points per flip, it is worth 0.875 points. You use it only on day 2, if the gap has reached $-2$.

Work back. On day 3 the gap goes to 0. So holding on day 2 is worth $|b|$: short if the gap is positive, long if it is negative. On day 1, the next move has zero mean, so the position does not matter in expectation, and the value is $E|b_2|$. From 10: $(14 + 6)/2 = 10$. From 2: $(6 + 2)/2 = 4$. From today: $(10 + 4)/2 = 7$.

The gain comes only from the path 6, 2, $-2$, of probability $\tfrac14$. Staying short loses 2 there, and flipping long makes 2. That is 4 more, and $\tfrac14 \times 4 = 1$. With the 0.5 cost, it is $\tfrac14 \times 3.5 = 0.875$.

This is a real part of the value of a basis trade: you hold an option on early exit and reversal, and that is one reason why people will hold a trade at a tight basis.

\iqwrong{Saying a basis that has to converge is worth just its gap. If you hold the right to reverse, the convergence pays more than the opening gap.}
\end{solution}

\begin{solution}{iq:iv2:etfs-futures-and-index-arbitrage:31}
You return about $-12.5\%$. The arbitrage does not close the premium because no one can create new shares. To short, you have to find a borrow, so you pay for one and you risk being squeezed. You would short small, sized for a further widening, and only with a catalyst that lets creation resume.

Your value is $1.05/1.20 = 0.875$. The underlying rose 5\%, and you still lost 12.5\%.

Normally a premium draws in creators, who sell the ETF against the basket. If creation is halted, the fund acts like a closed-end one, so its price can float away from value. The premium can widen before it closes, and the borrow may get expensive, so you can be right and still lose your position.

\iqwrong{Buying the fund for the underlying's move and not looking at the premium. You are buying the underlying and the premium, and the premium will not stay.}
\end{solution}

\begin{solution}{iq:iv2:etfs-futures-and-index-arbitrage:32}
Short interest above 100\% is possible because the same share can be lent, sold, then lent and sold again. The long holdings total 220 million: the 100 million outstanding plus 120 million from shorts. Market makers can create more shares to deliver or to meet demand to borrow. If the ETF closes, the main risk is that the fund's assets are realised, and shorts must buy back. Holders are paid through their brokers, with the shorts having to cover what they owe.

A lender holds a claim on a broker, and the buyer on the other side also counts as holding the share. So longs equal the shares outstanding plus the short interest. Each short is matched by an extra long.

The creation process gives ETFs a way of relieving shortage: an authorised participant can create shares, and deliver them to whoever needs them. So a high short interest in an ETF is often a hedging tool.

The risk is usually not on the assets, which are held for the fund. It is that an agent or a chain of lenders fails to deliver, or that in stress it takes time to unwind the chain. Investors can check how much is on loan.

\iqwrong{Saying that a short interest over 100\% means the fund has ``phantom shares'' and is about to collapse. The asset value is backed, and the extra longs are matched by shorts.}
\end{solution}

\begin{solution}{iq:iv2:etfs-futures-and-index-arbitrage:33}
The expected saving is 0.2 basis points of the fund, and the tracking risk is about 0.45 basis points of the fund. You could reasonably do either, and a fund that is judged on tracking error in the month might not.

The stock is 0.1\% of the fund, so the saved 2\% of impact is 0.2 basis points. The volatility over five days is $2\% \times \sqrt 5 \approx 4.47\%$, and on a 0.1\% weight, that is about 0.45 basis points. A position built evenly over the five days would risk about 0.30.

That is a fair Sharpe ratio in a single name. The problem is scale: the benefit is real only as long as few others do the same, and an index fund's mandate often makes it give up that gain. Many such changes over a year diversify the risk, and that argues for doing it.

\iqwrong{Comparing 3\% and 1\% without scaling to the fund. Both numbers on the fund's return are tiny, and what counts is how often this happens across all changes.}
\end{solution}

\begin{solution}{iq:iv2:etfs-futures-and-index-arbitrage:34}
The fund's value satisfies
\[
\frac{L_T}{L_0} = \left(\frac{S_T}{S_0}\right)^{\!L}\exp\left(-(L - 1)rT - \tfrac12L(L - 1)\sigma^2T - fT\right).
\]
For $L = 3$, $\sigma = 20\%$ and a flat index with zero rates, it is about 0.887. With a 10\% rise and rates of 4\%, it is about 1.090. An at-the-money one-year call on the fund, with zero rates, is worth about 23.6 per 100. Three such calls on the index are worth about 23.9.

The fund holds exposure $L$ and borrows $L - 1$ at $r$: $dL_t/L_t = L\,dS_t/S_t - (L - 1)r\,dt - f\,dt$. By Ito\index{Ito's lemma}, $d\ln S = dS/S - \tfrac12\sigma^2dt$ and $d\ln L_t = L\,dS/S - (L - 1)r\,dt - f\,dt - \tfrac12L^2\sigma^2dt$. Multiply the first by $L$ and subtract, and the $dS/S$ term disappears. So the fund is a deterministic function of today's index, and it has no other path dependence given a constant $\sigma$.

With numbers: $e^{-3 \times 0.04} = e^{-0.12} \approx 0.887$ at $r = 0$. With the rise and the rate, $1.331 \times e^{-0.08 - 0.12} \approx 1.090$.

The call: the fund has volatility $L\sigma = 60\%$, so its call is $100(2\Phi(0.3) - 1) \approx 23.6$. An index call is 7.97, and three of them make 23.9.

The two are close, but the payoffs are different. Three index calls pay linearly in the index, and the fund's call pays in its cube, with a scale factor. So the fund's call pays less if the index rises a little, and it pays more for a large rise.

\iqwrong{Assuming the formula holds if volatility changes. Then the exponent uses realised variance, and the fund depends on the path of volatility. So the fund's own value, and with it the payoff of an option on it, depends on realised variance, which the payoff of an index option does not.}
\end{solution}

\begin{solution}{iq:iv2:etfs-futures-and-index-arbitrage:35}
You expect about \$2 million on your \$100 million, or 2\%, with a standard deviation of about 4.5\%, which is about \$4.5 million. If others own \$400 million, the funds' net demand on the market is \$1.2 billion, the pressure is about 6.9\%, and half of that is left, so the profit falls to about 1.5\% before the extra competition for your entry. You would pass when many have the same position and the stock has already gone beyond the rule's price.

The peak is $2\% \times \sqrt{16} = 8\%$. Half was priced on the announcement, so 4\% is yet to come. Your entry costs $2\% \times \sqrt{1} = 2\%$: \$100 million is one day's volume. So the profit is $4\% - 2\% = 2\%$, or \$2 million.

The risk is the stock's move over the five days, $2\% \times \sqrt 5 \approx 4.47\%$, or \$4.47 million if you are unhedged. Then a hedge with its sector removes much of the market part. This is about 0.45 standard deviations per event, so it is a modest trade, even if it works.

With others holding \$400 million, net demand is \$1.2 billion, and $2\% \times \sqrt{12} \approx 6.93\%$. If half is priced, about 3.46\% remains, against your 2\% cost, which now rises because they are buying with you. That is about 1.46\%.

I would pass when the stock's move since the announcement is well above what this says, because arbitrageurs have already crowded it, or the funds' demand can be met by a large holder who is likely to sell at the close. And the risk of a reversal is greater if others all exit at the same time.

\iqwrong{Buying at the announcement, since ``the funds have to buy''. They do, and everyone knows, so the trade can be priced before you. And after the close, the price often falls back.}
\end{solution}
